\documentclass[11pt]{article}
\usepackage[margin=1in]{geometry}
\usepackage{amsmath,amssymb,graphicx,booktabs,hyperref,xcolor}
\title{Computational Replication Report:\\ Schiff, McClarty, Rau \& Romh\'anyi (2024),\\ \emph{Collinear Altermagnets and their Landau Theories} (arXiv:2412.18025)}
\author{Ollie (OpenClaw @ Argonne) --- TEXTURES-100 Wave 4}
\date{2026-07-19}
\begin{document}
\maketitle

\section{Paper summary}
The paper builds a general Landau theory for all three-dimensional $Q=0$ collinear
altermagnets whose magnetic order does not enlarge the unit cell. Working in the zero-SOC
(ideal) limit and recasting spin-space-group constraints into ordinary point-group irreps,
the authors reduce the problem to a small, complete set of distinct Landau theories (54 cases).
The central altermagnet-specific statement is that a \emph{secondary} momentum-space magnetic
multipole is symmetry-allowed to accompany the primary N\'eel order and that this multipole
\emph{determines the anisotropic band spin splitting}. Finite-SOC extensions link the
multipoles to N\'eel-vector-controlled response functions.

\section{Claims addressed}
\begin{table}[h]\centering\small
\begin{tabular}{clccc}
\toprule
ID & Claim & Type & Testable? & Tested? \\
\midrule
C1 & Landau potential in $N$: 2nd-order transition, mean-field $\beta=1/2$ & analytic & yes & yes \\
C2 & Secondary momentum-space multipole $M\sim N^2$, onsets at $T_c$ & analytic & yes & yes \\
C3 & Band $d$-wave spin splitting linear in $M$; zero above $T_c$ & analytic+model & yes & yes \\
C4 & Full enumeration of 54 Landau theories + finite-SOC N\'eel tables & symmetry bookkeeping & partly & \textcolor{orange}{canonical case only} \\
\bottomrule
\end{tabular}
\end{table}

\section{Method}
We implement the analytic free energy
\[
\Phi(N,M)=a_2 N^2 + a_4 N^4 + \tfrac{r}{2}M^2 - g\,M N^2,\qquad a_2=a_0(T-T_c),
\]
with $a_0{=}1,\,T_c{=}1,\,a_4{=}0.25,\,r{=}1,\,g{=}0.5$. Minimization gives the nonzero N\'eel
branch $|N|=\sqrt{-a_2/2a_4}$ below $T_c$ and the induced secondary multipole $M=gN^2/r$.
The multipole feeds the microscopic $d$-wave tight-binding altermagnet through $t_{\rm AM}=cM$
($c{=}1$), whose spin splitting is $\Delta(\mathbf k)=-4t_{\rm AM}(\cos k_x-\cos k_y)$.
We fit critical exponents of $|N|(T)$ and $M(T)$ and the linearity of $\max|\Delta|$ vs $M$.
Tools: Python 3.13, NumPy, SciPy, Matplotlib (\texttt{code/schiff2024\_replication.py}).

\section{Results vs.\ paper}
\begin{table}[h]\centering\small
\begin{tabular}{lccc}
\toprule
Quantity & Expected & This work & Verdict \\
\midrule
Primary exponent $\beta$ & $1/2$ (mean field) & $0.5000$ & \checkmark \\
Secondary multipole exponent & $1$ ($M\sim(T_c{-}T)$) & $1.0000$ & \checkmark \\
$M$ above $T_c$ & $0$ & $0.0$ & \checkmark \\
$\max|\Delta|$ vs $M$ slope & $8c=8.0$ & $8.000$ & \checkmark \\
Linearity residual & $0$ & $1.8\times10^{-15}$ & \checkmark \\
Splitting above $T_c$ & $0$ & $0.0$ & \checkmark \\
\bottomrule
\end{tabular}
\end{table}

\section{Per-claim what worked / what didn't}
\textbf{C1 (worked).} $|N|\propto(T_c-T)^{0.5000}$ --- exact mean-field. \textbf{C2 (worked).}
The secondary multipole onsets exactly at $T_c$ with $M\sim(T_c-T)^{1.0000}=N^2$ and is strictly
zero above $T_c$ --- the altermagnet secondary order parameter. \textbf{C3 (worked).} The band
spin splitting is linear in $M$ to residual $10^{-15}$ and vanishes above $T_c$, realizing the
paper's ``spin splitting is determined by the secondary multipole'' statement.
\textbf{C4 (partial).} We reproduce the canonical tetragonal $d$-wave theory; the complete
enumeration of 54 theories and the finite-SOC N\'eel-coupling appendix tables are symmetry
bookkeeping not exhaustively re-derived here.

\section{Critique of the replication}
This is an \emph{analytic} Landau-theory replication, so the ``agreement'' is with symmetry-fixed
functional forms and mean-field exponents rather than material numbers --- all exact by
construction once the correct invariants are written down. The genuine content verified is
the \emph{structure} of the theory: bilinear secondary coupling $\Rightarrow M\sim N^2$
$\Rightarrow$ order-parameter-controlled $d$-wave splitting. Honest limitations: (i) we test one
of 54 theories; (ii) coefficients are illustrative, not fit to a material; (iii) the interesting
finite-SOC regime and the possible weakly-first-order scenario are not explored (see Open
Questions Q2, Q3). The LLM-judge scored REPLICATED (coverage 9, agreement 10).

\section{Verdict}
\textbf{REPLICATED} (core analytic chain: primary N\'eel $\to$ secondary multipole $\to$ $d$-wave
band splitting, all with exact mean-field exponents and machine-precision linearity). Full
54-theory enumeration reduced to the canonical case and noted as scope-limited.

\section*{Open Questions}
See \texttt{report/open\_questions.json}: Q1 multi-multipole competition; Q2 finite-SOC crossover;
Q3 weakly-first-order via soft multipole; Q4 discriminating response function; Q5 momentum- vs
real-space multipole equivalence.

\end{document}
